\documentclass[11pt,a4paper]{article}
\usepackage[T1]{fontenc}
\usepackage[utf8]{inputenc}
\usepackage{lmodern,microtype}
\usepackage[a4paper,margin=23mm,headheight=15pt]{geometry}
\usepackage{amsmath,amssymb,amsthm,mathtools,booktabs,enumitem}
\usepackage{hyperref,graphicx,xcolor,fancyhdr}
\definecolor{accent}{HTML}{9B4635}
\definecolor{dark}{HTML}{1B2B2B}
\hypersetup{colorlinks=true,linkcolor=accent,urlcolor=accent,citecolor=accent}
\pagestyle{fancy}\fancyhf{}\lhead{\footnotesize\scshape Chapter 2: Time-Independent Schrödinger Equation}
\rhead{\footnotesize Azad}\cfoot{\thepage}
\numberwithin{equation}{section}
\newcommand{\dd}{\,\mathrm d}
\newcommand{\ii}{\mathrm i}
\newcommand{\ee}{\mathrm e}
\newcommand{\avg}[1]{\left\langle #1\right\rangle}
\newcommand{\kappaB}{\kappa}
\newcommand{\problem}[2]{\subsection*{\texorpdfstring{Problem #1: #2}{Problem #1: #2}}\addcontentsline{toc}{subsection}{Problem #1: #2}\label{prob:#1}}
\newcommand{\answer}[1]{\par\smallskip\noindent\textcolor{accent}{\textbf{Result.}}\quad #1\par\smallskip}
\newcommand{\qcheck}[1]{\par\smallskip\noindent\textit{Verification / interpretation.} #1\par\smallskip}
\newcommand{\basis}{\sqrt{2/a}}
\setlength{\parskip}{5pt}
\setlength{\parindent}{0pt}
\title{\vspace{-1.4cm}\textcolor{accent}{\large MATHEMATICAL PHYSICS / COURSE COMPANION}\\[12pt]
\textbf{\Huge Time-Independent\\[3pt] Schrödinger Equation}\\[13pt]
{\large Independent worked solutions and research-style commentary}\\[5pt]
{\normalsize Chapter 2, Problems 2.1--2.64}}
\author{Kumar Azad\\[-1pt]\small\url{https://krazad0.github.io/QM-CH1-SOLUTION/chapter2/}}
\date{November 2026}
\begin{document}
\maketitle
\begin{abstract}
This educational manuscript follows Chapter 2 of Griffiths and Schroeter,
\textit{Introduction to Quantum Mechanics}, third edition, and presents original
paraphrases and mathematical work for the chapter's 64 numbered problems.
The central themes are stationary-state boundary-value problems, orthonormal
eigenfunctions, superposition dynamics, harmonic-oscillator algebra, Fourier
wave packets, singular and finite-range potentials, scattering coefficients,
numerical spectral methods, and asymptotic analysis.
Long and computational problems include reproducible algorithms or reduced
equations; where a result requires numerical root finding, that requirement is
made explicit. This is a \textbf{coursework companion, not a new research result,
not a published arXiv preprint, and not an official solutions manual}.
\end{abstract}

\tableofcontents
\section{Notation and foundational identities}
The stationary eigenvalue equation and its associated time dependence are
\begin{align}
\hat H\psi_n&=E_n\psi_n,&
\hat H&=-\frac{\hbar^2}{2m}\frac{\dd^2}{\dd x^2}+V(x),&
\Psi_n(x,t)&=\psi_n(x)\ee^{-\ii E_nt/\hbar}.
\end{align}
When the eigenstates form a complete orthonormal discrete basis,
\begin{equation}
\Psi(x,t)=\sum_n c_n\psi_n(x)\ee^{-\ii E_nt/\hbar},\qquad
c_n=\int\psi_n^*(x)\Psi(x,0)\dd x,\qquad
\sum_n|c_n|^2=1.
\end{equation}
For the infinite well $0<x<a$, define
\begin{equation}
\phi_n(x)=\sqrt{\frac2a}\sin\left(\frac{n\pi x}{a}\right),\quad
E_n=n^2\hbar\Omega,\quad
\Omega=\frac{\pi^2\hbar}{2ma^2}.
\end{equation}
For the oscillator define $q=\sqrt{m\omega/\hbar}\,x$ and
$\ell=\sqrt{\hbar/(m\omega)}$.
Continuity of the wavefunction and derivative applies to finite jumps in a
regular potential; a Dirac delta of strength $\lambda$ instead imposes
$\psi'(x_0^+)-\psi'(x_0^-)=2m\lambda\psi(x_0)/\hbar^2$.
All statements concerning bound eigenstates assume appropriate self-adjoint
boundary conditions; formal plane waves require generalized normalization.

\section{Stationary states and the infinite square well (2.1--2.9)}
\problem{2.1}{Reality, parity, and stationary states}
\textit{Paraphrase.} Establish three properties of stationary solutions for real
one-dimensional, time-independent potentials.
For a separated solution $\Psi=\psi(x)\ee^{-\ii Et/\hbar}$, let
$E=E_R+\ii E_I$. Then
\begin{equation}
\|\Psi(t)\|^2=\ee^{2E_It/\hbar}\|\psi\|^2.
\end{equation}
A normalizable closed-system state preserves its norm, so $E_I=0$:
the energy is real. For real $V$, complex conjugation of
$-\hbar^2\psi''/(2m)+V\psi=E\psi$ shows that $\psi^*$ solves the same
equation. Hence $\Re\psi=(\psi+\psi^*)/2$ and
$\Im\psi=(\psi-\psi^*)/(2\ii)$ provide real solutions (at least one
nonzero) spanning the same real solution space. If $V(-x)=V(x)$, the
reflected function $\psi(-x)$ also solves the same equation.
Then
\begin{equation}
\psi_e(x)=\frac{\psi(x)+\psi(-x)}2,\qquad
\psi_o(x)=\frac{\psi(x)-\psi(-x)}2
\end{equation}
are even and odd eigenfunctions with the same energy.
\answer{Bound-state eigenfunctions can be chosen real and of definite
parity, and their energies are real.}
\check{In higher-dimensional or degenerate problems not every complex
eigenvector is a single parity eigenfunction; one chooses a suitable
basis within the eigenspace.}

\problem{2.2}{Why bound-state energy exceeds the potential minimum}
\textit{Paraphrase.} Prove a lower bound for any nontrivial normalizable eigenstate.
Multiply $\hat H\psi=E\psi$ by $\psi^*$ and integrate. Integration by
parts under boundary conditions giving vanishing surface terms gives
\begin{equation}
E\int|\psi|^2\dd x=
\frac{\hbar^2}{2m}\int|\psi'|^2\dd x+
\int V|\psi|^2\dd x
\ge V_{\min}\int|\psi|^2\dd x .
\end{equation}
For a normalized state,
\begin{equation}
E-V_{\min}=\frac{\hbar^2}{2m}\int|\psi'|^2\dd x
+\int(V-V_{\min})|\psi|^2\dd x.
\end{equation}
Both terms are nonnegative. Equality would require $\psi'$ to
vanish almost everywhere and the support to lie where $V=V_{\min}$;
a nonzero constant is not square-integrable on the real line.
\answer{$\boxed{E>V_{\min}}$ for a nontrivial normalizable bound
state on the infinite line under standard assumptions.}
Classically $E=T+V\ge V_{\min}$; the strict quantum inequality reflects
the impossibility of a normalizable state with everywhere vanishing
kinetic expectation.

\problem{2.3}{Excluding nonpositive well energies}
In the well $V=0$ for $0<x<a$ and $\psi(0)=\psi(a)=0$.
For $E=0$ the equation is $\psi''=0$, hence
$\psi=Ax+B$. The endpoint values yield $B=0$, then $A=0$.
For $E<0$, set $\kappa=\sqrt{-2mE}/\hbar>0$. Then
\begin{equation}
\psi=A\sinh(\kappa x)+B\cosh(\kappa x);
\quad \psi(0)=0\Rightarrow B=0;\quad
\psi(a)=0\Rightarrow A\sinh(\kappa a)=0.
\end{equation}
Since $\sinh(\kappa a)\ne0$, $A=0$ as well. Only the trivial state
remains. Thus $E_n=n^2\pi^2\hbar^2/(2ma^2)>0$.

\problem{2.4}{Position and momentum fluctuations in a box}
For $\phi_n=\sqrt{2/a}\sin(n\pi x/a)$,
\begin{align}
\avg{x}&=\frac2a\int_0^a x\sin^2(n\pi x/a)\dd x=\frac a2,\\
\avg{x^2}&=\frac2a\int_0^a x^2\sin^2(n\pi x/a)\dd x
=a^2\left(\frac13-\frac{1}{2n^2\pi^2}\right),\\
\avg p&=-\ii\hbar\int_0^a\phi_n\phi_n'\dd x
=-\frac{\ii\hbar}{2}[\phi_n^2]_0^a=0,\\
\avg{p^2}&=-\hbar^2\int_0^a\phi_n\phi_n''\dd x
=\frac{n^2\pi^2\hbar^2}{a^2}.
\end{align}
Therefore
\begin{equation}
\boxed{\sigma_x=a\sqrt{\frac1{12}-\frac1{2n^2\pi^2}},
\quad \sigma_p=\frac{n\pi\hbar}{a},\quad
\sigma_x\sigma_p=\hbar\sqrt{\frac{n^2\pi^2}{12}-\frac12}.}
\end{equation}
The expression increases with integer $n\ge1$, so the
ground state is closest to the Heisenberg bound:
$\sigma_x\sigma_p=\hbar\sqrt{\pi^2/12-1/2}\approx0.568\hbar$.
The strict inequality follows from the sine shape not being Gaussian.

\problem{2.5}{Equal two-state superposition in an infinite well}
Let $\Omega=\pi^2\hbar/(2ma^2)$ and
$\Psi(x,0)=A[\phi_1(x)+\phi_2(x)]$.
Orthonormality yields $1=|A|^2(1+1)$; choose $A=1/\sqrt2$.
The evolved state and its density are
\begin{align}
\Psi(x,t)&=\frac1{\sqrt2}
\left[\phi_1(x)\ee^{-\ii\Omega t}+
\phi_2(x)\ee^{-4\ii\Omega t}\right],\\
|\Psi|^2&=\frac12(\phi_1^2+\phi_2^2)
+\phi_1\phi_2\cos(3\Omega t).
\end{align}
The diagonal position elements are $a/2$, while
\begin{equation}
X_{12}=\int_0^a\phi_1(x)x\phi_2(x)\dd x
=\frac1a\int_0^a x[\cos(\pi x/a)-\cos(3\pi x/a)]\dd x
=-\frac{16a}{9\pi^2}.
\end{equation}
Hence
\begin{align}
\boxed{\avg{x}(t)=\frac a2-\frac{16a}{9\pi^2}\cos(3\Omega t)},\\
\boxed{\avg p(t)=m\frac{\dd\avg x}{\dd t}
=\frac{16ma\Omega}{3\pi^2}\sin(3\Omega t)}.
\end{align}
The oscillation amplitude is $16a/(9\pi^2)<a/2$ and angular
frequency $3\Omega$. Only $E_1=\hbar\Omega$ and
$E_2=4\hbar\Omega$ can be measured, with probability $1/2$
each; $\avg H=(E_1+E_2)/2=5\hbar\Omega/2$.

\problem{2.6}{Relative phase controls spatial oscillation}
Change the initial state to
$\Psi(x,0)=[\phi_1(x)+\ee^{\ii\varphi}\phi_2(x)]/\sqrt2$.
Time evolution gives
\begin{equation}
\Psi(x,t)=\frac{\phi_1\ee^{-\ii\Omega t}
+\ee^{\ii\varphi}\phi_2\ee^{-4\ii\Omega t}}{\sqrt2},
\end{equation}
so
\begin{align}
|\Psi|^2&=\frac12(\phi_1^2+\phi_2^2)
+\phi_1\phi_2\cos(3\Omega t-\varphi),\\
\boxed{\avg x(t)&=\frac a2-\frac{16a}{9\pi^2}
\cos(3\Omega t-\varphi).}
\end{align}
For $\varphi=\pi/2$ the oscillation is shifted by a quarter
period and $\avg x(0)=a/2$; for $\varphi=\pi$, the cosine term changes
sign and the initial spatial asymmetry reverses. The energy
probabilities remain $1/2$ and $1/2$. A global phase would cancel;
the \emph{relative} phase does not.

\problem{2.7}{Triangular initial state and spectral decomposition}
Let $\Psi(x,0)=A\min(x,a-x)$ on $[0,a]$, zero otherwise.
By symmetry,
\begin{equation}
1=2A^2\int_0^{a/2}x^2\dd x
=A^2\frac{a^3}{12}
\quad\Longrightarrow\quad A=\sqrt{\frac{12}{a^3}}.
\end{equation}
It has a peak $Aa/2$ at the midpoint. Its expansion coefficients are
\begin{align}
c_n&=\sqrt{\frac2a}A\left[
\int_0^{a/2}x\sin\frac{n\pi x}{a}\dd x+
\int_{a/2}^{a}(a-x)\sin\frac{n\pi x}{a}\dd x\right]\\
&=\boxed{\frac{4\sqrt6}{n^2\pi^2}
\sin\left(\frac{n\pi}{2}\right)}.
\end{align}
Thus all even terms vanish, and
\begin{equation}
\Psi(x,t)=\sum_{n=1,3,5,\ldots}c_n\phi_n(x)
\ee^{-\ii n^2\Omega t},\qquad
P(E_1)=|c_1|^2=\frac{96}{\pi^4}.
\end{equation}
Using $\sum_{n\ \mathrm{odd}\ge1}n^{-2}=\pi^2/8$,
\begin{align}
\avg H&=\sum_{n\ \mathrm{odd}}n^2\hbar\Omega|c_n|^2
=\hbar\Omega\frac{96}{\pi^4}\frac{\pi^2}{8}\\
&=\boxed{\frac{12\hbar\Omega}{\pi^2}
=\frac{6\hbar^2}{ma^2}}.
\end{align}
\check{The square-integrable jump in the first derivative does not
invalidate the spectral expansion; in Problem 2.39 it is treated
using distributions.}

\problem{2.8}{Energy measurement from a half-box constant state}
Here $\Psi(x,0)=A$ on $0<x<a/2$, and zero otherwise.
Normalization gives $|A|^2(a/2)=1$, so $A=\sqrt{2/a}$.
The ground-state overlap is
\begin{equation}
c_1=\sqrt{\frac2a}\sqrt{\frac2a}
\int_0^{a/2}\sin(\pi x/a)\dd x
=\frac{2}{a}\frac{a}{\pi}
=\frac2\pi.
\end{equation}
As the potential is time-independent, $|c_1|^2$ remains fixed:
\begin{equation}
\boxed{P\left(E=\frac{\pi^2\hbar^2}{2ma^2}\right)
=\frac4{\pi^2}\approx0.4053.}
\end{equation}

\problem{2.9}{Direct Hamiltonian expectation for a polynomial state}
For Example 2.2, $\Psi(x,0)=A x(a-x)$ with
$A=\sqrt{30/a^5}$. Inside the well,
\begin{equation}
\Psi''(x,0)=-2A,\qquad
\hat H\Psi=-\frac{\hbar^2}{2m}\Psi''
=\frac{\hbar^2 A}{m}.
\end{equation}
Then
\begin{align}
\avg H&=\int_0^a\Psi(x,0)\hat H\Psi(x,0)\dd x\\
&=\frac{\hbar^2A^2}{m}\int_0^a x(a-x)\dd x
=\frac{\hbar^2}{m}\frac{30}{a^5}\frac{a^3}{6}\\
&=\boxed{\frac{5\hbar^2}{ma^2}}.
\end{align}
This equals the spectral-sum result in Example 2.3.

\section{The harmonic oscillator (2.10--2.16)}
\problem{2.10}{Second excited harmonic oscillator state}
Introduce $q=\sqrt{m\omega/\hbar}\,x$ and the normalized
$\psi_0(x)=(m\omega/\pi\hbar)^{1/4}\ee^{-q^2/2}$.
Using
$a^\dagger=(q-\partial_q)/\sqrt2$ and
$\psi_2=(a^\dagger)^2\psi_0/\sqrt{2!}$,
\begin{equation}
a^\dagger\psi_0=\sqrt2\,q\psi_0,\quad
(a^\dagger)^2\psi_0=(2q^2-1)\psi_0,\quad
\boxed{\psi_2(x)=\frac{2q^2-1}{\sqrt2}\psi_0(x)}.
\end{equation}
The first three are $\psi_0=\psi_0(x)$ (even, no
nodes), $\psi_1=\sqrt2q\psi_0$ (odd, one node $q=0$), and
$\psi_2$ (even, nodes at $q=\pm1/\sqrt2$). The inner products
$\avg{\psi_0|\psi_1}$ and $\avg{\psi_1|\psi_2}$ vanish
by oddness. The nontrivial pair is
\begin{align}
\int_{\mathbb R}\psi_0\psi_2\,\dd x
&=\frac1{\sqrt2\sqrt\pi}
\int_{-\infty}^\infty(2q^2-1)\ee^{-q^2}\dd q\\
&=\frac{2(\sqrt\pi/2)-\sqrt\pi}{\sqrt{2\pi}}=0.
\end{align}
\qcheck{Normalization of $\psi_2$ follows from
$\avg{(2q^2-1)^2}_{\psi_0}/2
=[4(3/4)-4(1/2)+1]/2=1$.}

\problem{2.11}{Lowest two oscillator states: direct expectations}
For $\psi_0\propto\ee^{-q^2/2}$ and
$\psi_1=\sqrt2q\psi_0$, the densities are even, so
$\avg x=\avg p=0$ (the functions are real and the
momentum integral is a boundary term).
Use the Gaussian moments
\begin{equation}
\frac1{\sqrt\pi}\int q^2\ee^{-q^2}\dd q=\frac12,\qquad
\frac1{\sqrt\pi}\int q^4\ee^{-q^2}\dd q=\frac34.
\end{equation}
For the ground state, $\avg{q^2}_0=1/2$.
For the first excited state,
$\avg{q^2}_1=2\avg{q^4}_0=3/2$. Since
$x=\ell q$, with $\ell=\sqrt{\hbar/(m\omega)}$,
\begin{equation}
\boxed{\avg{x^2}_0=\frac{\hbar}{2m\omega},\qquad
\avg{x^2}_1=\frac{3\hbar}{2m\omega}}.
\end{equation}
Acting with $-\hbar^2\partial_x^2$ on $\psi_0$ gives
$\hat p^2\psi_0=(m\hbar\omega-m^2\omega^2x^2)\psi_0$;
on $\psi_1$ gives
$\hat p^2\psi_1=(3m\hbar\omega-m^2\omega^2x^2)\psi_1$.
Thus
\begin{equation}
\boxed{\avg{p^2}_0=\frac{m\hbar\omega}{2},\qquad
\avg{p^2}_1=\frac{3m\hbar\omega}{2}}.
\end{equation}
The uncertainty products are respectively
$\hbar/2$ and $3\hbar/2$. The potential and
kinetic energies in these states satisfy
\begin{equation}
\avg V=\frac12m\omega^2\avg{x^2}=\frac12 E_n,\qquad
\avg T=\frac{\avg{p^2}}{2m}=\frac12 E_n,
\end{equation}
giving totals $E_0=\hbar\omega/2$ and
$E_1=3\hbar\omega/2$ as required.

\problem{2.12}{General oscillator expectations using ladder algebra}
Use
\begin{equation}
x=\sqrt{\frac{\hbar}{2m\omega}}(a+a^\dagger),\qquad
p=\ii\sqrt{\frac{m\hbar\omega}{2}}(a^\dagger-a),
\end{equation}
together with $\langle n|a|n\rangle=0$,
$a^\dagger a|n\rangle=n|n\rangle$ and
$aa^\dagger|n\rangle=(n+1)|n\rangle$.
Terms with $a^2$ or $(a^\dagger)^2$ change $n$ by two
and therefore vanish in diagonal matrix elements.
Consequently
\begin{align}
\avg x_n&=\avg p_n=0,\\
\boxed{\avg{x^2}_n}&=\boxed{\frac{\hbar}{m\omega}
\left(n+\frac12\right)},\\
\boxed{\avg{p^2}_n}&=\boxed{m\hbar\omega
\left(n+\frac12\right)},\\
\boxed{\avg T_n}&=\boxed{\frac12\hbar\omega
\left(n+\frac12\right)},\\
\boxed{\sigma_x\sigma_p}&=\boxed{\hbar\left(n+\frac12\right)}.
\end{align}
This obeys the uncertainty relation for all $n\ge0$,
with equality only in the ground state.

\problem{2.13}{Dynamics of a $3:4$ oscillator superposition}
Let $\Psi(x,0)=A[3\psi_0(x)+4\psi_1(x)]$.
Orthonormality gives $A=1/5$.
Writing $E_n=(n+1/2)\hbar\omega$,
\begin{equation}
\Psi(x,t)=\frac35\psi_0\ee^{-\ii\omega t/2}
+\frac45\psi_1\ee^{-3\ii\omega t/2}.
\end{equation}
Because the stationary functions are real,
\begin{equation}
|\Psi|^2=\frac9{25}\psi_0^2+\frac{16}{25}\psi_1^2+
\frac{24}{25}\psi_0\psi_1\cos(\omega t).
\end{equation}
Replacing $\psi_1$ by $\psi_2$ would produce a density
cross-term oscillating at $2\omega$ instead. Since
$\langle0|x|1\rangle=\sqrt{\hbar/(2m\omega)}$,
\begin{align}
\boxed{\avg x(t)}&=\boxed{\frac{24}{25}
\sqrt{\frac{\hbar}{2m\omega}}\cos(\omega t)},\\
\boxed{\avg p(t)}&=\boxed{-\frac{24}{25}
\sqrt{\frac{m\hbar\omega}{2}}\sin(\omega t)}.
\end{align}
Differentiating the second equation gives
$\dd\avg p/\dd t=-m\omega^2\avg x$,
which is Ehrenfest's theorem.
Energy outcomes are $E_0$ and $E_1$ with
probabilities $9/25$ and $16/25$, and
\begin{equation}
\avg H=\frac9{25}\frac12\hbar\omega+
\frac{16}{25}\frac32\hbar\omega
=\boxed{\frac{57}{50}\hbar\omega}.
\end{equation}

\problem{2.14}{Ground-state probability outside classical turning points}
The ground-state energy is $\hbar\omega/2$.
Classical turning points satisfy
$m\omega^2x_t^2/2=E_0$, giving $x_t=\ell$
with $\ell=\sqrt{\hbar/(m\omega)}$.
For $q=x/\ell$, the normalized ground-state density is
$\rho_q=\ee^{-q^2}/\sqrt\pi$. Thus
\begin{equation}
P(|x|>x_t)=\frac{2}{\sqrt\pi}
\int_1^\infty \ee^{-q^2}\dd q
=\operatorname{erfc}(1)
=\boxed{0.157\quad\text{(three significant figures)}}.
\end{equation}
The nonzero tail is a direct signature of quantum
penetration into the classically forbidden region.

\problem{2.15}{Fifth and sixth Hermite polynomials}
The Hermite equation $H_n''-2qH_n'+2nH_n=0$
admits the recursion
$H_{n+1}=2qH_n-2nH_{n-1}$,
with $H_0=1$, $H_1=2q$.
Starting from $H_3=8q^3-12q$ and
$H_4=16q^4-48q^2+12$,
\begin{align}
H_5&=2q(16q^4-48q^2+12)-8(8q^3-12q)\\
&=\boxed{32q^5-160q^3+120q},\\
H_6&=2q(32q^5-160q^3+120q)-10(16q^4-48q^2+12)\\
&=\boxed{64q^6-480q^4+720q^2-120}.
\end{align}
The leading coefficients are $2^5$ and $2^6$ respectively.

\problem{2.16}{Rodrigues, recursion, differentiation, generating function}
For physicists' Hermite polynomials,
\begin{equation}
H_n(q)=(-1)^n\ee^{q^2}\frac{\dd^n}{\dd q^n}
\ee^{-q^2},\quad
H_{n+1}=2qH_n-2nH_{n-1}.
\end{equation}
Repeated differentiation in Rodrigues' expression gives
\begin{equation}
H_3=8q^3-12q,\qquad
H_4=16q^4-48q^2+12.
\end{equation}
The recurrence then gives the $H_5$ and $H_6$
calculated in Problem 2.15. Direct differentiation checks
\begin{align}
H_5'&=160q^4-480q^2+120=10H_4,\\
H_6'&=384q^5-1920q^3+1440q=12H_5.
\end{align}
Finally expand the generating function
\begin{equation}
\ee^{-z^2+2qz}
=1+2qz+(2q^2-1)z^2+
\left(\frac43q^3-2q\right)z^3+O(z^4).
\end{equation}
Reading off the coefficient of $z^n/n!$ gives
\begin{equation}
\boxed{H_1=2q,\quad H_2=4q^2-2,\quad H_3=8q^3-12q}.
\end{equation}
The three identities are mutually consistent and provide
independent calculations of the same polynomial family.

\section{Free particles and wave packets (2.17--2.21)}
\problem{2.17}{Traveling-wave and standing-wave bases}
Euler's identities imply
$\ee^{\ii kx}=\cos kx+\ii\sin kx$ and
$\ee^{-\ii kx}=\cos kx-\ii\sin kx$. Thus
\begin{equation}
A\ee^{\ii kx}+B\ee^{-\ii kx}
=(A+B)\cos kx+\ii(A-B)\sin kx.
\end{equation}
Comparing coefficients gives
\begin{equation}
\boxed{C=A+B,\qquad D=\ii(A-B)},\qquad
\boxed{A=\frac{C-\ii D}{2},\qquad
B=\frac{C+\ii D}{2}}.
\end{equation}
The two descriptions span the same two-dimensional solution
space for $E=\hbar^2k^2/(2m)$, but expose different physics:
exponentials carry directional flux, sines and cosines
describe standing waves.

\problem{2.18}{Probability current of free waves}
For a fixed positive $k$ and energy $E=\hbar^2k^2/(2m)$,
write
\begin{equation}
\Psi(x,t)=[A\ee^{\ii kx}+B\ee^{-\ii kx}]
\ee^{-\ii Et/\hbar}.
\end{equation}
The one-dimensional current is
$J=(\hbar/m)\operatorname{Im}(\Psi^*\partial_x\Psi)$.
The cross terms are real in $\Psi^*\partial_x\Psi$, so
\begin{equation}
\boxed{J=\frac{\hbar k}{m}(|A|^2-|B|^2)}.
\end{equation}
A single right-moving wave $A\ee^{\ii kx}$ has positive
current, a single left-moving wave $B\ee^{-\ii kx}$
negative current, and equal amplitudes produce a standing
wave with zero net current. An ideal plane wave is not
normalizable; $J$ describes flux per unit normalization
of the scattering amplitude.

\problem{2.19}{Complex Fourier series and the Fourier transform}
Start with an ordinary Fourier series on $[-a,a]$,
\begin{equation}
f(x)=\sum_{n=0}^{\infty}
\left[a_n\sin\frac{n\pi x}{a}+b_n\cos\frac{n\pi x}{a}\right].
\end{equation}
Using Euler's identity and taking $a_0=0$,
\begin{equation}
f(x)=\sum_{n=-\infty}^\infty c_n\ee^{\ii n\pi x/a},
\quad c_0=b_0,\quad
c_n=\frac{b_n-\ii a_n}{2},\quad
c_{-n}=\frac{b_n+\ii a_n}{2}\quad(n>0).
\end{equation}
Orthogonality of exponentials gives
\begin{equation}
\int_{-a}^{a}\ee^{\ii(m-n)\pi x/a}\dd x
=2a\delta_{mn}
\Rightarrow
\boxed{c_n=\frac1{2a}\int_{-a}^{a}
f(x)\ee^{-\ii n\pi x/a}\dd x}.
\end{equation}
Now define $k_n=n\pi/a$, $\Delta k=\pi/a$ and
$F(k_n)=\sqrt{2/\pi}\,a c_n$.
Then
\begin{equation}
f(x)=\frac1{\sqrt{2\pi}}\sum_{n=-\infty}^\infty
F(k_n)\ee^{\ii k_nx}\Delta k,\qquad
F(k_n)=\frac1{\sqrt{2\pi}}\int_{-a}^a
f(x)\ee^{-\ii k_nx}\dd x.
\end{equation}
In the limit $a\to\infty$, $\Delta k\to0$,
and for appropriate integrable functions these sums
and truncated integrals approach the Fourier transform pair
\begin{equation}
\boxed{F(k)=\frac1{\sqrt{2\pi}}\int_{\mathbb R}
f(x)\ee^{-\ii kx}\dd x,\qquad
f(x)=\frac1{\sqrt{2\pi}}\int_{\mathbb R}
F(k)\ee^{\ii kx}\dd k}.
\end{equation}
\qcheck{The finite-interval argument motivates the transform
rather than giving a rigorous existence theorem for all functions.
The Plancherel identity additionally preserves the $L^2$ norm.}

\problem{2.20}{Free evolution from a two-sided exponential}
For $\Psi(x,0)=A\ee^{-a|x|}$ with $a>0$,
\begin{equation}
1=2|A|^2\int_0^\infty\ee^{-2ax}\dd x
=\frac{|A|^2}{a}
\Rightarrow\quad \boxed{A=\sqrt a}.
\end{equation}
The Fourier coefficient is
\begin{align}
\phi(k)&=\frac{\sqrt a}{\sqrt{2\pi}}
\int_{-\infty}^\infty\ee^{-a|x|}
\ee^{-\ii kx}\dd x\\
&=\sqrt{\frac a{2\pi}}
\left[\frac1{a+\ii k}+\frac1{a-\ii k}\right]
=\boxed{\sqrt{\frac2\pi}\frac{a^{3/2}}{a^2+k^2}}.
\end{align}
Each momentum component evolves with frequency
$\omega_k=\hbar k^2/(2m)$. Therefore
\begin{equation}
\boxed{
\Psi(x,t)=\frac{a^{3/2}}{\pi}
\int_{-\infty}^\infty
\frac{\ee^{\ii kx-\ii\hbar k^2t/(2m)}}{a^2+k^2}\dd k}.
\end{equation}
When $a$ is large, the original state has width
$\sigma_x\sim1/a$ and a broad momentum distribution,
so it spreads rapidly; when $a$ is small the
wavefunction is very broad and its typical momentum
scale $\hbar a$ is small.
Neither extreme is a normalizable delta function
or uniform plane wave at the literal limiting parameter.

\problem{2.21}{Gaussian wave-packet spreading}
Take $\Psi(x,0)=A\ee^{-ax^2}$ with $a>0$.
Normalization gives $A=(2a/\pi)^{1/4}$.
Its Fourier transform is
\begin{equation}
\phi(k)=\frac{A}{\sqrt{2\pi}}
\int\ee^{-ax^2-\ii kx}\dd x
=\frac{A}{\sqrt{2a}}\ee^{-k^2/(4a)}.
\end{equation}
Completing the square in the inverse transform after
multiplication by $\ee^{-\ii\hbar k^2t/(2m)}$ yields
\begin{equation}
\boxed{\Psi(x,t)=
\left(\frac{2a}{\pi}\right)^{1/4}
\frac1{\sqrt{\gamma(t)}}
\exp\left[-\frac{ax^2}{\gamma(t)}\right]},
\qquad \gamma=1+\frac{2\ii\hbar at}{m}.
\end{equation}
Define $s=2\hbar at/m$ and
$w=a/(1+s^2)$. Since
$\operatorname{Re}(1/\gamma)=1/(1+s^2)$,
\begin{equation}
\boxed{|\Psi|^2=\sqrt{\frac{2w}{\pi}}
\ee^{-2wx^2}}.
\end{equation}
The state remains Gaussian but its width grows.
Parity gives $\avg{x}=\avg p=0$.
The second moments and uncertainties are
\begin{align}
\avg{x^2}&=\frac{1+s^2}{4a},&
\avg{p^2}&=a\hbar^2,\\
\boxed{\sigma_x}&=\boxed{\frac{\sqrt{1+s^2}}{2\sqrt a}},&
\boxed{\sigma_p}&=\boxed{\hbar\sqrt a},\\
\boxed{\sigma_x\sigma_p}&=\boxed{\frac{\hbar}{2}\sqrt{1+s^2}}.
\end{align}
The momentum uncertainty remains constant, whereas
position uncertainty increases due to free dispersion.
The lower bound is saturated at $t=0$ alone.
\qcheck{A wavepacket's density can evolve even though each
Fourier mode has a fixed energy; interference between
different energies generates the spreading.}

\section{Delta potentials (2.22--2.28)}
\problem{2.22}{Three delta-function integrals}
The sifting identity is
$\int_I f(x)\delta(x-x_0)\dd x=f(x_0)$
if $x_0$ lies in the interior of $I$, and zero if it
lies outside. Consequently
\begin{align}
\int_{-3}^{1}(x^3-3x^2+2x-1)\delta(x+2)\dd x
&=(-2)^3-3(-2)^2+2(-2)-1=\boxed{-25},\\
\int_0^\infty[\cos(3x)+2]\delta(x-\pi)\dd x
&=\cos(3\pi)+2=\boxed{1},\\
\int_{-1}^{1}\ee^{|x|+3}\delta(x-2)\dd x
&=\boxed{0}.
\end{align}
The third integral vanishes because the delta's support
lies outside the domain, regardless of the integrand
evaluated at its center.

\problem{2.23}{Scaling of the delta and distributional derivative of a step}
For nonzero real $c$ and a smooth test function $f$,
substitute $u=cx$, taking account of orientation if
$c<0$:
\begin{equation}
\int_{-\infty}^{\infty}f(x)\delta(cx)\dd x
=\frac1{|c|}\int f(u/c)\delta(u)\dd u
=\frac{f(0)}{|c|}.
\end{equation}
Therefore $\boxed{\delta(cx)=\delta(x)/|c|}$ as
distributions. For the Heaviside step $\theta$,
integration by parts against a compactly supported
smooth test function gives
\begin{equation}
\langle\theta',f\rangle
=-\int_0^\infty f'(x)\dd x=f(0)
=\langle\delta,f\rangle,
\end{equation}
and hence $\boxed{\theta'(x)=\delta(x)}$
in the distributional sense. The conventional
value $\theta(0)=1/2$ has no effect on this result.

\problem{2.24}{Uncertainty product for a delta-bound state}
For $V=-\alpha\delta(x)$ with $\alpha>0$, define
$\kappa=m\alpha/\hbar^2$. The normalized bound state is
\begin{equation}
\psi(x)=\sqrt\kappa\,\ee^{-\kappa|x|},
\qquad E=-\frac{\hbar^2\kappa^2}{2m}.
\end{equation}
Symmetry gives $\avg x=\avg p=0$ and
\begin{equation}
\avg{x^2}=2\kappa\int_0^\infty
x^2\ee^{-2\kappa x}\dd x=\frac1{2\kappa^2}.
\end{equation}
The first derivative is
$\psi'=-\kappa\,\operatorname{sgn}(x)\psi$,
whereas the second derivative contains a delta:
\begin{equation}
\psi''=\kappa^2\psi-2\kappa\psi(0)\delta(x).
\end{equation}
Therefore, using $\avg{p^2}=-\hbar^2\int\psi^*\psi''\dd x$,
\begin{equation}
\avg{p^2}=-\hbar^2\kappa^2+
2\hbar^2\kappa|\psi(0)|^2
=\hbar^2\kappa^2.
\end{equation}
Consequently
\begin{equation}
\boxed{\sigma_x=\frac1{\sqrt2\kappa},\quad
\sigma_p=\hbar\kappa,\quad
\sigma_x\sigma_p=\frac{\hbar}{\sqrt2}>
\frac{\hbar}{2}}.
\end{equation}
\check{The delta term is essential; ignoring the derivative
jump would incorrectly produce a negative $\avg{p^2}$.}

\problem{2.25}{Bound-scattering orthogonality in the delta well}
Let $\psi_b$ be the normalized delta-well bound state
with $E_b<0$ and let $\psi_k$ be a generalized scattering
eigenfunction with $E_k>0$. Both satisfy the same matching
condition at $x=0$:
$\psi'(0^+)-\psi'(0^-)=-2m\alpha\psi(0)/\hbar^2$.
Multiply the equations $\hat H\psi_b=E_b\psi_b$
and $\hat H\psi_k=E_k\psi_k$ by $\psi_k^*$
and $\psi_b$, subtract and integrate on the two
half-lines. The Wronskian boundary contributions at
$0$ cancel because both states obey the same jump,
and the terms at infinity vanish because
$\psi_b,\psi_b'\to0$ exponentially and $\psi_k$ is bounded.
Thus
\begin{equation}
(E_k-E_b)\int_{-\infty}^\infty\psi_b(x)\psi_k^*(x)\dd x=0.
\end{equation}
Since $E_k\ne E_b$,
\begin{equation}
\boxed{\int\psi_b(x)\psi_k^*(x)\dd x=0}.
\end{equation}
\check{The scattering state is not individually in
$L^2(\mathbb R)$, but this mixed overlap is convergent.}

\problem{2.26}{Fourier transform of a delta distribution}
For the symmetric convention
$\widehat f(k)=(2\pi)^{-1/2}
\int f(x)\ee^{-\ii kx}\dd x$,
\begin{equation}
\boxed{\widehat\delta(k)=\frac1{\sqrt{2\pi}}}.
\end{equation}
Formally applying the inverse transform gives
\begin{equation}
\boxed{\delta(x)=\frac1{2\pi}
\int_{-\infty}^\infty\ee^{\ii kx}\dd k}.
\end{equation}
This is an identity of distributions, not an
ordinary convergent improper integral.
For instance, a Gaussian regulator gives
\begin{align}
\frac1{2\pi}\int\ee^{-\epsilon k^2}
\ee^{\ii kx}\dd k
&=\frac1{\sqrt{4\pi\epsilon}}
\ee^{-x^2/(4\epsilon)}
\longrightarrow\delta(x)
\quad(\epsilon\downarrow0).
\end{align}
Against any smooth, rapidly decaying test function,
the right-hand side converges to its value at $x=0$.

\problem{2.27}{Bound states of two attractive delta centers}
For $V=-\alpha[\delta(x-a)+\delta(x+a)]$
with $\kappa_0=m\alpha/\hbar^2$ and $E<0$,
set $\kappa=\sqrt{-2mE}/\hbar>0$. By parity, a
normalized bound state has the form
\begin{align}
\psi_e(x)&=\begin{cases}
C\cosh(\kappa x),&|x|<a,\\
C\cosh(\kappa a)\ee^{-\kappa(|x|-a)},&|x|\ge a,
\end{cases}\\
\psi_o(x)&=\begin{cases}
C\sinh(\kappa x),&|x|<a,\\
C\,\operatorname{sgn}(x)\sinh(\kappa a)
\ee^{-\kappa(|x|-a)},&|x|\ge a.
\end{cases}
\end{align}
Continuity at $a$ is automatic. Imposing the derivative
jump $\psi'(a^+)-\psi'(a^-)=-2\kappa_0\psi(a)$ gives
\begin{equation}
\boxed{\kappa_e=\kappa_0(1+\ee^{-2\kappa_e a}),\qquad
\kappa_o=\kappa_0(1-\ee^{-2\kappa_o a})}.
\end{equation}
The even state always exists; a positive odd root exists
only for $2\kappa_0a>1$. For the specified strengths:
\begin{center}
\begin{tabular}{lcc}
\toprule
$\alpha$&$\kappa_e a$&$\kappa_o a$\\
\midrule
$\hbar^2/(ma)$&$1.108858$&$0.796812$\\
$\hbar^2/(4ma)$&$0.369418$&none\\
\bottomrule
\end{tabular}
\end{center}
Each energy is $E=-\hbar^2(\kappa a)^2/(2ma^2)$.
As $a\to0$, the even state approaches a single
delta of strength $2\alpha$, with
$E_e\to-2m\alpha^2/\hbar^2$; the odd state disappears.
As $a\to\infty$, the two bound levels converge
to the isolated-well energy $-m\alpha^2/(2\hbar^2)$.

\problem{2.28}{Transmission through the double-delta potential}
For $k=\sqrt{2mE}/\hbar>0$ and
$\beta=m\alpha/(\hbar^2k)$, write a local plane-wave
solution as $A\ee^{\ii kx}+B\ee^{-\ii kx}$ to
the left and $F\ee^{\ii kx}+G\ee^{-\ii kx}$
to the right. Matching at an attractive delta at
$x=x_0$ yields
\begin{equation}
\begin{pmatrix}F\\G\end{pmatrix}=
\underbrace{\begin{pmatrix}
1+\ii\beta&\ii\beta\ee^{-2\ii kx_0}\\
-\ii\beta\ee^{2\ii kx_0}&1-\ii\beta
\end{pmatrix}}_{M(x_0)}
\begin{pmatrix}A\\B\end{pmatrix}.
\end{equation}
For centers at $-a$ and $a$, $M=M(a)M(-a)$.
Because $\det M=1$, left incidence ($G=0$) gives
$F/A=1/M_{22}$, where
\begin{equation}
M_{22}=(1-\ii\beta)^2+\beta^2\ee^{4\ii ka}.
\end{equation}
Thus
\begin{equation}
\boxed{T(k)=
\frac1{\left|(1-\ii\beta)^2+
\beta^2\ee^{4\ii ka}\right|^2}}.
\end{equation}
This is real, nonnegative, and no larger than one
for a real potential. As $a\to0$,
$T\to[1+4\beta^2]^{-1}$, precisely the result for
a single attractive delta of strength $2\alpha$.

\section{Finite wells, barriers and steps (2.29--2.35)}
\problem{2.29}{Odd bound states of a finite square well}
For $V(x)=-V_0$ on $|x|<a$ and $0$ outside,
bound levels satisfy $-V_0<E<0$. Define
\begin{equation}
l=\frac{\sqrt{2m(E+V_0)}}{\hbar},\quad
\kappa=\frac{\sqrt{-2mE}}{\hbar},\quad
z=la,\quad z_0=\frac{a\sqrt{2mV_0}}{\hbar}.
\end{equation}
An odd eigenfunction is $D\sin(lx)$ inside and
$F\operatorname{sgn}(x)\ee^{-\kappa|x|}$ outside.
At $x=a$,
\begin{equation}
F\ee^{-\kappa a}=D\sin(la),\qquad
-\kappa F\ee^{-\kappa a}=lD\cos(la).
\end{equation}
Dividing yields $l\cot(la)=-\kappa$, or
\begin{equation}
\boxed{-z\cot z=\sqrt{z_0^2-z^2}},\qquad 0<z<z_0.
\end{equation}
This equation can be solved by intersections of
$-z\cot z$ with $\sqrt{z_0^2-z^2}$ on the
intervals where the left side is positive.
For a deep well ($z_0\gg1$), roots approach
$z=n\pi/2$ for even $n$, corresponding to
odd-parity states of the centered infinite well.
As the well weakens the first odd state appears
only when $z_0>\pi/2$; there is no odd bound
state for arbitrarily weak $V_0$.
The even ground state, by contrast, always exists.

\problem{2.30}{Normalizing a finite-well even bound state}
The even solution is
\begin{equation}
\psi(x)=
\begin{cases}
D\cos(lx),&|x|<a,\\
F\ee^{-\kappa|x|},&|x|\ge a.
\end{cases}
\end{equation}
Continuity at $x=a$ gives
$F\ee^{-\kappa a}=D\cos(la)$,
or $F=D\cos(la)\ee^{\kappa a}$.
Hence
\begin{align}
1&=2D^2\int_0^a\cos^2(lx)\dd x+
2F^2\int_a^\infty\ee^{-2\kappa x}\dd x\\
&=D^2\left[a+\frac{\sin(2la)}{2l}
+\frac{\cos^2(la)}{\kappa}\right].
\end{align}
Choose the overall phase so $D>0$:
\begin{equation}
\boxed{D=\left[a+\frac{\sin(2la)}{2l}
+\frac{\cos^2(la)}{\kappa}\right]^{-1/2}},
\quad
\boxed{F=D\cos(la)\ee^{\kappa a}}.
\end{equation}
The even-state eigenvalue relation $\kappa=l\tan(la)$
from the textbook must be solved before evaluating
these constants numerically.

\problem{2.31}{The delta well as the zero-width square-well limit}
Let the square well of half-width $a$ have depth
$V_0=\alpha/(2a)$, so its integrated strength
is $-2aV_0=-\alpha$. As $a\to0$,
\begin{equation}
z_0^2=\frac{2mV_0a^2}{\hbar^2}
=\frac{m\alpha a}{\hbar^2}\longrightarrow0.
\end{equation}
For the even bound level $\kappa=l\tan(la)$ and
$la=z\to0$, so
\begin{equation}
\kappa\sim l^2a
=\frac{2m(E+V_0)a}{\hbar^2}
\longrightarrow\frac{m\alpha}{\hbar^2}.
\end{equation}
Therefore
\begin{equation}
\boxed{E\longrightarrow
-\frac{\hbar^2\kappa^2}{2m}
=-\frac{m\alpha^2}{2\hbar^2}},
\end{equation}
the single attractive-delta bound energy.
For positive incident energy,
the square-well transmission formula is
\begin{equation}
T^{-1}=1+\frac{V_0^2}{4E(E+V_0)}
\sin^2\left(\frac{2a\sqrt{2m(E+V_0)}}{\hbar}\right).
\end{equation}
Using $\sin z\sim z$ and $2aV_0=\alpha$,
\begin{equation}
T^{-1}\to1+\frac{m\alpha^2}{2\hbar^2 E}
=1+\left(\frac{m\alpha}{\hbar^2k}\right)^2,
\quad k=\sqrt{2mE}/\hbar.
\end{equation}
The apparent infinite depth is therefore
a weak-well limit in the finite-square-well
dimensionless parameter $z_0$.

\problem{2.32}{Finite square-well scattering}
For $E>0$ define
$k=\sqrt{2mE}/\hbar$ outside and
$l=\sqrt{2m(E+V_0)}/\hbar$ inside
the well $-a<x<a$.
Write the left and right solutions
\begin{equation}
\psi_L=A\ee^{\ii kx}+B\ee^{-\ii kx},\qquad
\psi_R=F\ee^{\ii kx}
\end{equation}
and the interior as
$\psi_I=C\sin(lx)+D\cos(lx)$.
The matching equations for both $\psi$ and
$\psi'$ at $\pm a$ give
\begin{equation}
\frac{A}{F}
=\ee^{2\ii ka}
\left[\cos(2la)-\frac{\ii}{2}
\left(\frac{k}{l}+\frac{l}{k}\right)\sin(2la)\right].
\end{equation}
Taking the squared modulus and using
$k^2-l^2=-2mV_0/\hbar^2$,
\begin{align}
\frac{|A|^2}{|F|^2}
&=\cos^2(2la)+\frac14
\left(\frac{k}{l}+\frac{l}{k}\right)^2\sin^2(2la)\\
&=1+\frac{(l^2-k^2)^2}{4k^2l^2}\sin^2(2la).
\end{align}
Since exterior velocities agree, $T=|F/A|^2$:
\begin{equation}
\boxed{T=
\left[1+\frac{V_0^2}{4E(E+V_0)}
\sin^2\left(\frac{2a\sqrt{2m(E+V_0)}}{\hbar}\right)
\right]^{-1}}.
\end{equation}
Resonances occur when $2la=n\pi$, making $T=1$.

\problem{2.33}{Transmission across a rectangular barrier}
For a barrier $V_0>0$ on $|x|<a$,
its full width is $L=2a$. Repeating the matching
of Problem 2.32 with interior wave number
$q=\sqrt{2m(E-V_0)}/\hbar$ gives for $E>V_0$
\begin{equation}
\boxed{T^{-1}=1+
\frac{V_0^2}{4E(E-V_0)}
\sin^2\left(\frac{L\sqrt{2m(E-V_0)}}{\hbar}\right)}.
\end{equation}
For $E<V_0$ put
$q=\ii\kappa$, $\kappa=\sqrt{2m(V_0-E)}/\hbar$,
and use $\sin(\ii y)=\ii\sinh y$. This yields
\begin{equation}
\boxed{T^{-1}=1+
\frac{V_0^2}{4E(V_0-E)}\sinh^2(\kappa L)}.
\end{equation}
The $E=V_0$ case is the continuous limit:
$\sin(qL)/q\to L$, leading to
\begin{equation}
\boxed{T(E=V_0)=\left[
1+\frac{mV_0 L^2}{2\hbar^2}\right]^{-1}}.
\end{equation}
For $E<V_0$, $T$ is exponentially suppressed at
large $\kappa L$, but is not zero: quantum tunneling.

\problem{2.34}{Reflection and transmission at an upward step}
Let $V=0$ for $x<0$ and $V=V_0$ for $x>0$.
The incident wave has $k=\sqrt{2mE}/\hbar$.
For $0<E<V_0$ the right-hand solution is
$F\ee^{-\kappa x}$ with
$\kappa=\sqrt{2m(V_0-E)}/\hbar$. Continuity of
$\psi$ and $\psi'$ at zero yields
\begin{equation}
\frac BA=\frac{k-\ii\kappa}{k+\ii\kappa},
\qquad \boxed{R=1,\quad T=0}.
\end{equation}
An evanescent tail exists for $x>0$ but carries no
transmitted flux. For $E>V_0$ let
$q=\sqrt{2m(E-V_0)}/\hbar$ and
$\psi_R=F\ee^{\ii qx}$.
Matching $A+B=F$, $k(A-B)=qF$ yields
\begin{equation}
\boxed{\frac BA=\frac{k-q}{k+q},
\quad \frac FA=\frac{2k}{k+q}}.
\end{equation}
The current of a wave $C\ee^{\ii kx}$ is
$\hbar k|C|^2/m$, so
\begin{equation}
\boxed{R=\frac{(k-q)^2}{(k+q)^2}},\qquad
\boxed{T=\frac qk\frac{|F|^2}{|A|^2}
=\frac{4kq}{(k+q)^2}},
\end{equation}
and $R+T=1$. The factor $q/k$ is essential
because the particle speed changes across the step.

\problem{2.35}{Quantum reflection from a downward step}
For $V=0$ on $x<0$ and $V=-V_0$ on $x>0$,
the transmitted wave number is
$q=\sqrt{2m(E+V_0)}/\hbar$.
Matching as before gives
\begin{equation}
R=\left(\frac{k-q}{k+q}\right)^2
=\left(\frac{\sqrt E-\sqrt{E+V_0}}
{\sqrt E+\sqrt{E+V_0}}\right)^2 .
\end{equation}
At $E=V_0/3$, the ratio $q/k=\sqrt{1+3}=2$,
so $\boxed{R=1/9}$ and $\boxed{T=8/9}$.
This also applies to the simplified example of
a $4$ MeV neutron approaching a $12$ MeV nuclear
potential drop: the probability of
\emph{entering} the interior is $8/9$ in
this one-dimensional step model.
\qcheck{Transmission through the nuclear surface
is not identical to irreversible nuclear absorption;
that additional process requires a more complete
interaction model. A falling car at a cliff also does
not experience a sudden change of gravitational
potential energy at the edge; the actual potential
depends on height and changes during the fall.}

\section{Further mathematical problems (2.36--2.45)}
\problem{2.36}{Centered infinite square well}
For $V=0$ on $-a<x<a$ and $V=\infty$ elsewhere,
the boundary conditions are $\psi(-a)=\psi(a)=0$.
Even solutions have the form $\cos(kx)$ and satisfy
$\cos(ka)=0$, giving $ka=(j+1/2)\pi$.
Odd solutions have the form $\sin(kx)$ and satisfy
$\sin(ka)=0$, giving $ka=j\pi$.
A combined ordering is
\begin{equation}
\boxed{E_n=\frac{n^2\pi^2\hbar^2}{8ma^2}},\qquad
n=1,2,\ldots,
\end{equation}
with
\begin{equation}
\boxed{\psi_n(x)=
\begin{cases}
a^{-1/2}\cos(n\pi x/2a),&n\text{ odd},\\
a^{-1/2}\sin(n\pi x/2a),&n\text{ even}.
\end{cases}}
\end{equation}
Signs can be changed without physical effect.
Each integral $\int_{-a}^a|\psi_n|^2\dd x=1$.
These are the familiar well states for a box of
width $2a$, shifted to the origin. The first three
are an even ground state with no internal node,
an odd first excited state with one internal node,
and an even second excited state with two nodes.

\problem{2.37}{Cubic sine initial wavefunction}
On $0<x<a$, let $\Psi(x,0)=A\sin^3(\pi x/a)$.
Use $\sin^3\theta=(3\sin\theta-\sin3\theta)/4$ and
$\int_0^\pi\sin^6\theta\dd\theta=5\pi/16$.
Normalization yields
\begin{equation}
1=A^2\frac{5a}{16}
\quad\Rightarrow\quad \boxed{A=\frac4{\sqrt{5a}}}.
\end{equation}
The eigenstate decomposition is
\begin{equation}
\Psi(x,0)=\frac3{\sqrt{10}}\phi_1(x)
-\frac1{\sqrt{10}}\phi_3(x).
\end{equation}
Therefore
\begin{equation}
\boxed{\Psi(x,t)=\frac3{\sqrt{10}}\phi_1(x)\ee^{-\ii\Omega t}
-\frac1{\sqrt{10}}\phi_3(x)\ee^{-9\ii\Omega t}}.
\end{equation}
Both $\phi_1,\phi_3$ are even about $x=a/2$.
The density remains symmetric about the midpoint,
so $\boxed{\avg x(t)=a/2}$ for every $t$.
Energy probabilities are $9/10$ and $1/10$:
\begin{equation}
\boxed{\avg H=\frac9{10}E_1+\frac1{10}E_3
=\frac95\hbar\Omega}.
\end{equation}
\qcheck{The density is not stationary, despite a
constant mean position; interference changes its shape.}

\problem{2.38}{Quantum revivals and classical crossing times}
Any normalized well state has expansion
$\Psi(x,t)=\sum_n c_n\phi_n(x)
\ee^{-\ii n^2\Omega t}$.
Choose $T_q=2\pi/\Omega$; then
$\ee^{-\ii n^2\Omega T_q}=\ee^{-2\pi\ii n^2}=1$
for every integer $n$. Hence
\begin{equation}
\boxed{\Psi(x,T_q)=\Psi(x,0)},\qquad
\boxed{T_q=\frac{4ma^2}{\pi\hbar}}.
\end{equation}
A classical particle traverses the box and returns
after a round-trip distance $2a$ at speed
$v=\sqrt{2E/m}$:
\begin{equation}
\boxed{T_{\mathrm{cl}}=\frac{2a}{v}
=2a\sqrt{\frac{m}{2E}}}.
\end{equation}
Equating the two times gives
\begin{equation}
\boxed{E=\frac{\pi^2\hbar^2}{8ma^2}=\frac{E_1}{4}}.
\end{equation}
This is below the lowest allowed energy of the
infinite well. The two concepts of recurrence
are therefore not equal for any stationary
state in the well.

\problem{2.39}{Second derivative of a state with a cusp}
From Problem 2.7,
\begin{equation}
\Psi(x,0)=A
\begin{cases}x,&0<x<a/2,\\a-x,&a/2<x<a,\end{cases}
\quad A=\sqrt{12/a^3}.
\end{equation}
Inside the open interval $(0,a)$, the first
derivative is the step-function expression
\begin{equation}
\boxed{\Psi'(x,0)=A[1-2\theta(x-a/2)]}.
\end{equation}
Using $\theta'=\delta$,
\begin{equation}
\boxed{\Psi''(x,0)=-2A\delta(x-a/2)}.
\end{equation}
This is a distributional derivative, not an ordinary
second derivative at the cusp. The quadratic energy
form gives
\begin{align}
\avg H&=-\frac{\hbar^2}{2m}\int_0^a
\Psi(x,0)\Psi''(x,0)\dd x\\
&=\frac{\hbar^2}{m}A\Psi(a/2,0)
=\frac{\hbar^2}{m}A^2\frac a2
=\boxed{\frac{6\hbar^2}{ma^2}}.
\end{align}
It agrees with Problem 2.7. A rigorous
alternative is the kinetic-energy quadratic form
$\hbar^2\int_0^a|\Psi'|^2/(2m)$.

\problem{2.40}{A quadratic polynomial times a Gaussian}
Set $q=\sqrt{m\omega/\hbar}\,x$.
The initial state is proportional to
$(1-2q)^2\ee^{-q^2/2}$.
Since $H_0=1$, $H_1=2q$, and
$H_2=4q^2-2$,
\begin{equation}
(1-2q)^2=3H_0-2H_1+H_2.
\end{equation}
Using normalized oscillator functions
$\psi_n=\alpha H_n(q)\ee^{-q^2/2}/\sqrt{2^nn!}$
where $\alpha=(m\omega/\pi\hbar)^{1/4}$,
\begin{equation}
\Psi(x,0)=A_{\mathrm{norm}}
\left[3\psi_0-2\sqrt2\psi_1+2\sqrt2\psi_2\right].
\end{equation}
The squared norm of the bracket is
$9+8+8=25$, so $A_{\mathrm{norm}}=1/5$.
The prefactor multiplying the \emph{unnormalized}
polynomial-Gaussian is $A=\alpha/5$.
Thus
\begin{equation}
\boxed{(c_0,c_1,c_2)=
\left(\frac35,-\frac{2\sqrt2}5,
\frac{2\sqrt2}5\right)}.
\end{equation}
Energy probabilities are $9/25$, $8/25$, $8/25$ and
\begin{equation}
\avg H=
\left[\frac9{25}\frac12+\frac8{25}\frac32+
\frac8{25}\frac52\right]\hbar\omega
=\boxed{\frac{73}{50}\hbar\omega}.
\end{equation}
At $t=\pi/\omega$, relative phases are
$\ee^{-\ii n\pi}=(-1)^n$.
The odd $n=1$ coefficient changes sign while
$n=0,2$ do not, so the polynomial becomes
$(1+2q)^2$ up to a global phase.
The smallest positive time is
$\boxed{T=\pi/\omega}$.

\problem{2.41}{The half harmonic oscillator}
For $x>0$, $V=m\omega^2x^2/2$; on $x\le0$
there is an infinite wall. Thus
$\psi(0)=0$ and $\psi(x)$ must decay as
$x\to+\infty$. Extending a half-line solution
oddly to the entire line produces an
odd full-oscillator eigenfunction.
Only $n=1,3,5,\dots$ are permitted.
Reindex $n=2j+1$, $j=0,1,\dots$:
\begin{equation}
\boxed{E_j=\left(2j+\frac32\right)\hbar\omega}.
\end{equation}
The normalized half-line wavefunctions are
$\boxed{\psi_j^{\rm half}(x)=\sqrt2\,
\psi_{2j+1}^{\rm full}(x)}$ for $x>0$.
The factor $\sqrt2$ compensates for
discarding the negative half of a normalized
odd-state density.

\problem{2.42}{Traveling Gaussian packet: envelope and phase speeds}
For $\Psi(x,0)=A\ee^{-ax^2}\ee^{\ii\ell x}$,
normalization still gives $A=(2a/\pi)^{1/4}$.
The initial Fourier distribution is shifted:
$\phi(k)\propto\ee^{-(k-\ell)^2/(4a)}$.
In the free-particle Fourier integral set
$u=k-\ell$ and complete the Gaussian square.
With $\gamma=1+2\ii\hbar at/m$,
\begin{equation}
\boxed{\Psi(x,t)=
\left(\frac{2a}{\pi}\right)^{1/4}
\gamma^{-1/2}
\exp\left[-\frac{a(x-v_gt)^2}{\gamma}\right]
\exp\left[\ii\ell\left(x-\frac{\hbar\ell t}{2m}\right)\right]},
\quad v_g=\frac{\hbar\ell}{m}.
\end{equation}
The density is a translated spreading Gaussian:
\begin{equation}
|\Psi|^2=
\sqrt{\frac{2a}{\pi|\gamma|^2}}\,
\exp\left[-\frac{2a(x-v_gt)^2}{|\gamma|^2}\right].
\end{equation}
Thus the envelope travels at
$v_g=\hbar\ell/m$; its central phase has
$\omega_\ell=\hbar\ell^2/(2m)$, so
the phase speed is
$v_p=\omega_\ell/\ell=\hbar\ell/(2m)=v_g/2$
for $\ell\ne0$. The packet satisfies
$\avg x(t)=v_gt$, $\avg p=\hbar\ell$.
\qcheck{The local phase of the dispersing Gaussian
envelope also evolves, so this $v_p$ refers to
its plane-wave carrier, not to every local
phase contour of the full packet.}

\problem{2.43}{Infinite well with a central delta barrier}
Take $V(x)=\alpha\delta(x)$ for $|x|<a$,
infinite outside, with $\alpha>0$.
Odd states vanish at $x=0$ and satisfy
the delta jump automatically. They remain
\begin{equation}
\psi_o(x)\propto\sin(j\pi x/a),\qquad
\boxed{E_{o,j}=\frac{j^2\pi^2\hbar^2}{2ma^2}}
\quad(j=1,2,\ldots).
\end{equation}
An even state on the right half may be written
$\psi_e(x)=C\sin[k(a-x)]$ for $0<x<a$,
reflected evenly for $x<0$.
At the origin the jump is
$2\psi_e'(0^+)=2m\alpha\psi_e(0)/\hbar^2$,
giving
\begin{equation}
-2Ck\cos(ka)=\frac{2m\alpha}{\hbar^2}C\sin(ka)
\quad\Rightarrow\quad
\boxed{-ka\cot(ka)=\frac{m\alpha a}{\hbar^2}}.
\end{equation}
At $\alpha=0$, the roots are $ka=(j+1/2)\pi$.
At $\alpha\to\infty$, the even roots approach
$ka=j\pi$ from below and become degenerate
with the unaffected odd states. A repulsive
central barrier raises the even-state energies,
consistent with the variational principle.

\problem{2.44}{No degenerate bound levels on a connected line}
Let $\psi_1,\psi_2$ solve the same
one-dimensional stationary equation at energy $E$:
\begin{equation}
-\frac{\hbar^2}{2m}\psi_j''+V\psi_j=E\psi_j.
\end{equation}
Multiply the first by $\psi_2$, the second by
$\psi_1$ and subtract:
$\psi_2\psi_1''-\psi_1\psi_2''=0$.
This is the derivative of the Wronskian
$W=\psi_2\psi_1'-\psi_1\psi_2'$.
Hence $W$ is constant. For ordinary bound states
$\psi_j,\psi_j'\to0$ at an outer boundary, so
$W=0$ everywhere. Wherever $\psi_1\ne0$,
\begin{equation}
\left(\frac{\psi_2}{\psi_1}\right)'=
\frac{\psi_1\psi_2'-\psi_1'\psi_2}{\psi_1^2}=0.
\end{equation}
Thus $\psi_2=C\psi_1$ and continuity extends the
relation through zeros. The solutions are not
linearly independent. Therefore each isolated
bound energy is nondegenerate on a connected
one-dimensional line under these assumptions.

\problem{2.45}{Node interlacing for different bound energies}
Take real eigenfunctions $\psi_m,\psi_n$ with
$E_n>E_m$ and define
\begin{equation}
W(x)=\psi_m'(x)\psi_n(x)-\psi_m(x)\psi_n'(x).
\end{equation}
Use $\psi_j''=(2m/\hbar^2)(V-E_j)\psi_j$
to obtain
\begin{equation}
\boxed{W'(x)=\frac{2m}{\hbar^2}
(E_n-E_m)\psi_m(x)\psi_n(x)}.
\end{equation}
Between successive nodes $x_1<x_2$ of $\psi_m$,
integrate to get
\begin{equation}
\boxed{\psi_m'(x_2)\psi_n(x_2)-
\psi_m'(x_1)\psi_n(x_1)=
\frac{2m(E_n-E_m)}{\hbar^2}
\int_{x_1}^{x_2}\psi_m\psi_n\dd x}.
\end{equation}
Suppose $\psi_n$ has no zero in $(x_1,x_2)$.
Then it has a definite sign, as does $\psi_m$.
If $\psi_m>0$ in this interval, then
$\psi_m'(x_1)>0$ and $\psi_m'(x_2)<0$;
the left side has the \emph{opposite} sign
to $\psi_n$ while the right side has the
\emph{same} sign. This is a contradiction.
The argument with $\psi_m<0$ is analogous.
Thus $\psi_n$ has a node between adjacent
nodes of $\psi_m$. With Sturm--Liouville
boundary conditions this interlacing yields
the increasing node count with energy.

\section{Dynamics, symmetries, and scattering matrices (2.46--2.54)}
\problem{2.46}{Stationary quantum states on a ring}
A ring of circumference $L$ imposes periodicity
$\psi(x+L)=\psi(x)$. Free solutions are
$\ee^{\ii kx}$; periodicity requires
$k_nL=2\pi n$ for integer $n$.
Normalization over one circumference gives
\begin{equation}
\boxed{\psi_n(x)=\frac1{\sqrt L}\ee^{2\pi\ii nx/L}},
\qquad
\boxed{E_n=\frac{2\pi^2\hbar^2n^2}{mL^2}},
\quad n\in\mathbb Z.
\end{equation}
For each positive $n$, the independent
$\psi_n$ and $\psi_{-n}$ carry opposite currents
and have equal energy. The $n=0$ state is
nondegenerate. The Wronskian theorem from
Problem 2.44 assumes wavefunctions decay
at spatial infinity on a connected line;
a circle is periodic rather than a line
with vanishing tails, so its constant
Wronskian need not be zero.

\problem{2.47}{Qualitative behavior of a double finite well}
At zero separation, the two attractive wells merge
into one wider region, and the lowest wavefunction
is nodeless and broadly distributed; the first excited
function has one antisymmetric node. At intermediate
separation the lowest pair becomes a symmetric and
antisymmetric combination of wavefunctions localized
near the left and right wells. At very large separation,
the overlap between the two centers becomes tiny:
the even and odd combinations have nearly the
same energy and become effectively independent.
The ground energy tends toward the isolated-well
ground energy from below, whereas the antisymmetric
partner tends toward it from above.
The splitting shrinks as the barrier separation
increases. In this model the electron favors closer
nuclei because the symmetric bonding state's
energy falls as the wells overlap. Real molecules
also involve internuclear repulsion, electron--electron
effects and three-dimensional geometry, so the
simple potential is only an analogy.
\check{This is deliberately a qualitative answer:
the textbook specifies no calculation for this problem.}

\problem{2.48}{Bound states beside a hard wall}
The potential is infinite for $x<0$, equals
$-V_0$ for $0<x<a$, and is zero for $x>a$,
where $V_0=32\hbar^2/(ma^2)$. For bound
states $-V_0<E<0$, set
\begin{equation}
z=ka=\frac{a\sqrt{2m(E+V_0)}}{\hbar},
\quad q=\kappa a=\frac{a\sqrt{-2mE}}{\hbar},
\quad z^2+q^2=64.
\end{equation}
The wall forces $\psi=C\sin(kx)$ inside;
outside, $\psi=C\sin z\,\ee^{-\kappa(x-a)}$.
The derivative matching gives
\begin{equation}
\boxed{-z\cot z=\sqrt{64-z^2}},\quad0<z<8.
\end{equation}
The three roots of this equation occur at
\begin{equation}
z\simeq2.785902,\quad5.521446,\quad7.957321,
\end{equation}
so there are three bound states. Their energies are
$E=-\hbar^2q^2/(2ma^2)$ with
$q=\sqrt{64-z^2}$.
For the highest state, $q\simeq0.825248$.
Integrating its density on either side of $a$,
\begin{align}
I_{\mathrm{in}}&=C^2a\left(
\frac12-\frac{\sin(2z)}{4z}\right),\\
I_{\mathrm{out}}&=C^2a\frac{\sin^2z}{2q},\\
P_{\mathrm{out}}&=\frac{I_{\mathrm{out}}}{I_{\mathrm{in}}+
I_{\mathrm{out}}}
=\boxed{0.542041\ \text{(approximately)}}.
\end{align}
\check{Despite $E<0$, the shallowest bound state
has more probability outside the finite well
than inside it.}

\problem{2.49}{Exact coherent oscillator wave packet}
Write $\ell^2=\hbar/(m\omega)$ and
$X(t)=x_0\cos\omega t$. The exact packet
given in the textbook can be written
\begin{align}
\Psi(x,t)&=(\pi\ell^2)^{-1/4}
\exp\Bigg[-\frac1{2\ell^2}
\left(x^2+\frac{x_0^2}2
(1+\ee^{-2\ii\omega t})
-2x_0x\ee^{-\ii\omega t}\right)
-\frac{\ii\omega t}{2}\Bigg].
\end{align}
To check it, denote its exponent by $S$ and use
$\partial_x\Psi=S_x\Psi$,
$\partial_x^2\Psi=(S_{xx}+S_x^2)\Psi$,
$\partial_t\Psi=S_t\Psi$. Inserting
\begin{align}
S_x&=-x/\ell^2+x_0\ee^{-\ii\omega t}/\ell^2,\\
S_{xx}&=-1/\ell^2,\\
S_t&=\frac{\ii\omega x_0^2}{\ell^2}\ee^{-2\ii\omega t}
-\frac{\ii\omega x_0 x}{\ell^2}\ee^{-\ii\omega t}
-\frac{\ii\omega}{2}
\end{align}
into $i\hbar S_t=
-\hbar^2(S_{xx}+S_x^2)/(2m)+m\omega^2x^2/2$
verifies the Schrödinger equation term by term.
Taking the real part of $S$,
\begin{equation}
\boxed{|\Psi(x,t)|^2=
\frac1{\sqrt{\pi}\ell}
\exp\left[-\frac{(x-x_0\cos\omega t)^2}{\ell^2}\right]}.
\end{equation}
The width remains constant:
$\sigma_x=\ell/\sqrt2$.
Gaussian integration gives
$\boxed{\avg x=x_0\cos\omega t}$ and
$\boxed{\avg p=-m\omega x_0\sin\omega t}$.
Then $\dd\avg p/\dd t=-m\omega^2\avg x$,
precisely Ehrenfest's theorem.

\problem{2.50}{A translating attractive delta well}
Let $V=-\alpha\delta(x-vt)$ and
$\kappa=m\alpha/\hbar^2$.
Start with the stationary bound state
$\psi_0(u)=\sqrt\kappa\ee^{-\kappa|u|}$
at energy $E_0=-m\alpha^2/(2\hbar^2)$.
A Galilean boost of velocity $v$ gives
\begin{equation}
\boxed{\Psi(x,t)=
\sqrt\kappa\ee^{-\kappa|x-vt|}
\exp\left[\frac{\ii}{\hbar}
\left(mvx-\left(E_0+\frac12mv^2\right)t\right)\right]}.
\end{equation}
For $u=x-vt$, differentiation of the phase
produces a kinetic-energy shift $mv^2/2$
and cross terms that cancel the translation
term $-v\partial_u\psi_0$. The remaining
equation is exactly the stationary delta-well
equation, including the distributional cusp.
Therefore this is a normalized exact solution.
Its instantaneous Hamiltonian expectation is
\begin{equation}
\boxed{\avg H=E_0+\frac12mv^2}.
\end{equation}
The center of the density moves at $vt$ while
the internal binding energy remains $E_0$.

\problem{2.51}{Gaussian free fall in a constant force}
Let $\Psi_0(y,t)$ be a freely evolving Gaussian
with $\avg y_0=0$ and zero mean momentum.
Define $y=x+gt^2/2$ and
\begin{equation}
\boxed{\Psi(x,t)=
\Psi_0\left(x+\frac12gt^2,t\right)
\exp\left[-\frac{\ii mgt}{\hbar}
\left(x+\frac16gt^2\right)\right]}.
\end{equation}
If $\Phi$ denotes the phase, its derivatives are
$\partial_x\Phi=-mgt/\hbar$ and
$\partial_t\Phi=-mg(x+gt^2/2)/\hbar$.
The chain rule gives
$\partial_t\Psi_0(y,t)=
(\partial_t\Psi_0)_y+gt\,\partial_y\Psi_0$;
the extra convective derivative cancels the
cross term from differentiating the phase
in the kinetic operator. The remaining term
is the required gravitational potential $mgx$.
Since the phase has unit modulus, the density
is simply translated:
\begin{equation}
|\Psi(x,t)|^2=
|\Psi_0(x+gt^2/2,t)|^2,\qquad
\boxed{\avg x(t)=-\frac12gt^2}.
\end{equation}
The center follows classical free fall although
the Gaussian simultaneously spreads.

\problem{2.52}{A reflectionless hyperbolic-secant well}
Let
\begin{equation}
V(x)=-\frac{\hbar^2a^2}{m}\operatorname{sech}^2(ax).
\end{equation}
The potential is even, approaches zero at both
infinities, and reaches $-\hbar^2a^2/m$ at zero.
For $\psi_0=A\operatorname{sech}(ax)$,
the identity
\begin{equation}
\frac{\psi_0''}{\psi_0}=
a^2[1-2\operatorname{sech}^2(ax)]
\end{equation}
shows
$\hat H\psi_0=-\hbar^2a^2\psi_0/(2m)$.
As $\int_{-\infty}^{\infty}
\operatorname{sech}^2(ax)\dd x=2/a$,
\begin{equation}
\boxed{\psi_0=\sqrt{\frac a2}\operatorname{sech}(ax),
\qquad E_0=-\frac{\hbar^2a^2}{2m}}.
\end{equation}
For $E=\hbar^2k^2/(2m)>0$,
take the exact scattering solution
\begin{equation}
\psi_k(x)=A\frac{\ii k-a\tanh(ax)}{\ii k+a}
\ee^{\ii kx}.
\end{equation}
Using $(\tanh ax)'=a\operatorname{sech}^2(ax)$
and $(\operatorname{sech}^2 ax)'=
-2a\operatorname{sech}^2(ax)\tanh(ax)$
in $\psi_k''$ verifies $\hat H\psi_k=E\psi_k$.
At $x\to-\infty$, $\tanh(ax)\to-1$,
so $\psi_k\to A\ee^{\ii kx}$, with no
left-moving reflected component. At
$x\to+\infty$,
\begin{equation}
\psi_k\to A
\frac{\ii k-a}{\ii k+a}\ee^{\ii kx}.
\end{equation}
The amplitude ratio has magnitude one, hence
\begin{equation}
\boxed{R=0,\qquad T=1}
\end{equation}
at every positive incident energy: an exactly
reflectionless potential.

\problem{2.53}{Scattering matrices for symmetric localized wells}
For a real localized potential, let $A,G$ be
the incoming amplitudes from left and right,
and $B,F$ the outgoing amplitudes:
\begin{equation}
\begin{pmatrix}B\\F\end{pmatrix}
=S\begin{pmatrix}A\\G\end{pmatrix}.
\end{equation}
For the attractive delta well at $x=0$ set
$\beta=m\alpha/(\hbar^2k)$.
The matching conditions give amplitudes
$r=\ii\beta/(1-\ii\beta)$ and
$t=1/(1-\ii\beta)$. Left--right symmetry gives
\begin{equation}
\boxed{S_{\delta}=
\frac1{1-\ii\beta}
\begin{pmatrix}\ii\beta&1\\1&\ii\beta\end{pmatrix}}.
\end{equation}
Indeed $|r|^2+|t|^2=1$ and
$r^*t+t^*r=0$, so $S^\dagger S=I$.
For a finite symmetric well of half-width $a$,
write $k=\sqrt{2mE}/\hbar$,
$l=\sqrt{2m(E+V_0)}/\hbar$,
and
$D=\cos(2la)-\ii(k^2+l^2)\sin(2la)/(2kl)$.
Solving the four continuity equations yields
\begin{equation}
\boxed{t=\frac{\ee^{-2\ii ka}}{D},\qquad
r=\frac{\ii(l^2-k^2)\sin(2la)\ee^{-2\ii ka}}{2klD}},
\qquad
\boxed{S_{\rm well}=\begin{pmatrix}r&t\\t&r\end{pmatrix}}.
\end{equation}
The probabilities $|t|^2$ and $|r|^2$
satisfy $T+R=1$; parity symmetry makes
the left and right reflection/transmission
probabilities equal.

\problem{2.54}{Transfer matrices, composition, and double deltas}
A transfer matrix relates right-hand amplitudes
to left-hand amplitudes:
\begin{equation}
\begin{pmatrix}F\\G\end{pmatrix}
=M\begin{pmatrix}A\\B\end{pmatrix}.
\end{equation}
Starting with
$B=S_{11}A+S_{12}G$ and
$F=S_{21}A+S_{22}G$, solve for $G$ and $F$:
\begin{equation}
\boxed{M=\frac1{S_{12}}
\begin{pmatrix}-\det S&S_{22}\\-S_{11}&1\end{pmatrix}}.
\end{equation}
Conversely, solving $G=M_{21}A+M_{22}B$
for $B$, then substituting into $F$ gives
\begin{equation}
\boxed{S=\frac1{M_{22}}
\begin{pmatrix}-M_{21}&1\\\det M&M_{12}\end{pmatrix}}.
\end{equation}
Hence
\begin{equation}
R_l=\left|\frac{M_{21}}{M_{22}}\right|^2,\quad
T_l=\left|\frac{\det M}{M_{22}}\right|^2,\quad
R_r=\left|\frac{M_{12}}{M_{22}}\right|^2,\quad
T_r=\frac1{|M_{22}|^2}.
\end{equation}
For a real localized potential $\det M=1$.
Two isolated potentials traversed from left to
right act consecutively, so
$\boxed{M_{\rm total}=M_2M_1}$.
For $V=-\alpha\delta(x-x_0)$, the
jump condition gives
\begin{equation}
\boxed{M(x_0)=
\begin{pmatrix}
1+\ii\beta&\ii\beta\ee^{-2\ii kx_0}\\
-\ii\beta\ee^{2\ii kx_0}&1-\ii\beta
\end{pmatrix}},
\qquad\beta=\frac{m\alpha}{\hbar^2k}.
\end{equation}
The two-delta transfer matrix is
$M(a)M(-a)$. Its lower-right element is
$M_{22}=(1-\ii\beta)^2+\beta^2\ee^{4\ii ka}$,
so
\begin{equation}
\boxed{T_{\delta\delta}=
\frac1{\left|(1-\ii\beta)^2+\beta^2\ee^{4\ii ka}\right|^2}},
\end{equation}
agreeing with Problem 2.28.

\section{Numerical and advanced problems (2.55--2.64)}
\problem{2.55}{Numerical shooting for the oscillator ground state}
With $q=\sqrt{m\omega/\hbar}\,x$ and
$K=2E/(\hbar\omega)$, the dimensionless equation is
\begin{equation}
u''(q)=(q^2-K)u(q).
\end{equation}
For an even ground state use $u(0)=1,u'(0)=0$,
integrate to a large cutoff $q=b$, and vary $K$
until the growing asymptotic component changes sign
and the tail becomes decaying. A robust alternative
is to match logarithmic derivatives of inward and
outward solutions, reducing sensitivity to the
exponentially growing solution.
The exact result provides a stringent five-digit check:
\begin{equation}
\boxed{K_0=1.0000,\qquad
E_0=0.50000\,\hbar\omega}.
\end{equation}
\check{A plot that appears to decay on a finite interval
does not by itself prove an eigenvalue. Verify node
count and cutoff convergence.}

\problem{2.56}{First three excited oscillator energies by shooting}
Solve the same equation as Problem 2.55.
For odd parity impose $u(0)=0,u'(0)=1$;
for even parity impose $u(0)=1,u'(0)=0$.
In increasing order the first three excited
solutions have one, two, and three nodes and
five-significant-digit dimensionless eigenvalues
\begin{equation}
\boxed{K_1=3.0000,\quad K_2=5.0000,\quad K_3=7.0000}.
\end{equation}
Since $E=(\hbar\omega/2)K$,
\begin{equation}
\boxed{E_1=1.5000\,\hbar\omega,\quad
E_2=2.5000\,\hbar\omega,\quad
E_3=3.5000\,\hbar\omega}.
\end{equation}
Alternating parity and rejecting growing tails
distinguishes the successive bound states.

\problem{2.57}{Numerical shooting for four infinite-well energies}
Scale the box by $y=x/a\in[0,1]$ and set
$\varepsilon=2ma^2E/\hbar^2$. The equation becomes
\begin{equation}
-u''(y)=\varepsilon u(y),\qquad
u(0)=0,\quad u'(0)=1,\quad u(1)=0.
\end{equation}
With trial $\varepsilon>0$,
$u(y)=\sin(\sqrt{\varepsilon}\,y)/\sqrt\varepsilon$.
The shooting residual is
$u(1)=\sin\sqrt{\varepsilon}/\sqrt\varepsilon$,
whose roots give
\begin{equation}
\boxed{\varepsilon_1=9.8696,\quad
\varepsilon_2=39.478,\quad
\varepsilon_3=88.826,\quad
\varepsilon_4=157.91}.
\end{equation}
The exact benchmark is $\varepsilon_n=n^2\pi^2$.
The corresponding eigenfunctions have $n-1$
interior nodes.

\problem{2.58}{Simple estimate of metallic cohesion}
An isolated electron in a width-$a$ box has
ground energy $e_1=\pi^2\hbar^2/(2ma^2)$.
For $N$ separate atoms with one electron each,
$E_{\rm sep}=Ne_1$.
Place the $N$ electrons (temporarily ignoring spin)
in a single box of width $Na$, one in each
of its lowest $N$ levels. The total is
\begin{align}
E_{\rm metal}
&=\sum_{n=1}^N\frac{n^2\pi^2\hbar^2}{2mN^2a^2}\\
&=e_1\frac{N(N+1)(2N+1)}{6N^2}
=e_1\frac{(N+1)(2N+1)}{6N}.
\end{align}
Thus the energy required to separate the atoms is
$E_{\rm coh}=E_{\rm sep}-E_{\rm metal}$.
For large $N$,
\begin{equation}
\boxed{\frac{E_{\rm coh}}{N}\longrightarrow
\frac23 e_1=\frac{\pi^2\hbar^2}{3ma^2}}.
\end{equation}
At $a=4$ \AA{} with the electron mass,
\begin{equation}
\boxed{E_{\rm coh}/N\simeq1.57\ \mathrm{eV}}.
\end{equation}
This toy model omits spin, three-dimensional
bands and ion-core structure, but captures
the lowering of kinetic energy when electrons
delocalize across a larger volume.

\problem{2.59}{Quantum bouncing ball: Airy eigenstates}
For a particle above a reflecting floor, let
$V=mgx$ when $x>0$ and infinite below.
Define
\begin{equation}
a_g=\left(\frac{2m^2g}{\hbar^2}\right)^{1/3},
\quad z=a_gx,\quad
\lambda=\frac{a_gE}{mg}
=\frac{2mE}{\hbar^2a_g^2}.
\end{equation}
The Schrödinger equation becomes
$-y''(z)+zy(z)=\lambda y(z)$,
whose general solution is
$C\operatorname{Ai}(z-\lambda)+
D\operatorname{Bi}(z-\lambda)$.
The $\operatorname{Bi}$ term diverges as
$z\to+\infty$, so $D=0$. The floor enforces
$\operatorname{Ai}(-\lambda_n)=0$.
The first and tenth positive Airy-zero magnitudes
are
\begin{equation}
\boxed{\lambda_1=2.338107410,\qquad
\lambda_{10}=12.828776753}.
\end{equation}
The Airy normalization identity gives
\begin{equation}
\int_0^\infty\operatorname{Ai}(z-\lambda_n)^2\dd z
=\operatorname{Ai}'(-\lambda_n)^2,
\qquad
\boxed{\psi_n(x)=
\frac{\sqrt{a_g}\operatorname{Ai}(a_gx-\lambda_n)}
{|\operatorname{Ai}'(-\lambda_n)|}}.
\end{equation}
In units $a_g=1$, the normalizing
prefactors are approximately $1.426105$
and $0.936510$ for the two states.
Their orthogonality follows from the
self-adjoint Sturm--Liouville equation,
and can also be verified numerically.
Virial and Airy integral identities give
\begin{equation}
\avg z=\frac23\lambda,\qquad
\avg{z^2}=\frac8{15}\lambda^2,\qquad
\avg{p^2}=\hbar^2a_g^2\,\frac{\lambda}{3},
\quad\avg p=0.
\end{equation}
Therefore
\begin{equation}
\boxed{\sigma_x=\frac{2\lambda}{\sqrt{45}\,a_g},
\quad \sigma_p=\hbar a_g\sqrt{\frac{\lambda}{3}},
\quad \frac{\sigma_x\sigma_p}{\hbar}
=\frac{2\lambda^{3/2}}{\sqrt{135}}}.
\end{equation}
For $\lambda_1$ this is $0.61540$, and
for $\lambda_{10}$ it is $7.90935$,
both above $1/2$.
Classically, the particle has a height cutoff
$z=\lambda$ and residence-time density
\begin{equation}
\boxed{\rho_C(z)=\frac1{2\sqrt{\lambda(\lambda-z)}}},
\quad0<z<\lambda.
\end{equation}
This diverges integrably at the classical
turning point, while the quantum density
$|\psi_n|^2$ oscillates and extends
beyond it; the coarse-grained high-$n$
profile approaches the classical envelope.
\check{The companion Python numerics script
plots these comparisons in the scaled coordinates.}

\problem{2.60}{Inverse-square attraction and the boundary-condition anomaly}
For $V=-\alpha/x^2$ on $x>0$, the constants
$m,\hbar,\alpha$ provide no independent length
scale: $\beta=2m\alpha/\hbar^2$ is dimensionless.
Hence there is no nonzero energy uniquely fixed
by dimensional analysis alone.
For $E=-\hbar^2\kappa^2/(2m)$,
\begin{equation}
\psi''+\frac{\beta}{x^2}\psi=\kappa^2\psi.
\end{equation}
If $\psi(x)$ solves this for $\kappa$,
$\psi(sx)$ solves it for $s\kappa$, giving
the formal scaling $E\mapsto s^2E$.
Set $g=\sqrt{\beta-1/4}$, assumed real
($\beta>1/4$). Let $\psi(x)=\sqrt x\,f(\kappa x)$;
the equation for $f(z)$ is
\begin{equation}
z^2f''+zf'-(z^2-g^2)f=0.
\end{equation}
The decaying solution is the modified Bessel
function $K_{\ii g}$:
\begin{equation}
\boxed{\psi_\kappa(x)=C\sqrt x\,K_{\ii g}(\kappa x)}.
\end{equation}
Its norm uses the identity
$\int_0^\infty z\,K_{\ii g}(z)^2\dd z
=\pi g/(2\sinh\pi g)$, so
\begin{equation}
\boxed{C=\kappa\sqrt{\frac{2\sinh(\pi g)}{\pi g}}}.
\end{equation}
As $x\to0$, the wavefunction behaves as
$\sqrt x$ times an oscillatory function of $\ln x$,
with infinitely many nodes approaching zero.
\textit{Important mathematical qualification:}
square integrability alone does not choose
a self-adjoint Hamiltonian at the singular origin.
A fixed self-adjoint extension supplies a
boundary phase/length scale, so formal
continuous rescaling need not preserve its
domain. This is the scale anomaly of the
attractive inverse-square problem.
If an infinite wall is instead placed at
$x=\ell>0$, the new condition is
\begin{equation}
\boxed{K_{\ii g}(\kappa\ell)=0}.
\end{equation}
Its largest positive $\kappa$ root
gives the most negative bound level.
Writing $z_*(g)=\kappa_0\ell$,
\begin{equation}
\boxed{E_0=-\frac{\alpha}{\ell^2}
f(\beta)},\qquad
\boxed{f(\beta)=\frac{z_*(\sqrt{\beta-1/4})^2}{\beta}}.
\end{equation}
This displays explicitly how a length parameter
restores a well-defined energy scale.

\problem{2.61}{Finite-difference Hamiltonian for a box}
Discretize the interval $[0,a]$ into $N+1$
equal segments of length $\Delta x=a/(N+1)$.
With Dirichlet endpoints $\psi_0=\psi_{N+1}=0$,
the $N\times N$ tridiagonal matrix is
\begin{equation}
\boxed{H_N=\frac{\hbar^2}{2m(\Delta x)^2}
\begin{pmatrix}
2&-1&0&\cdots\\
-1&2&-1&\cdots\\
0&-1&2&\ddots\\
\vdots&\vdots&\ddots&\ddots
\end{pmatrix}}.
\end{equation}
In units $E_*= \hbar^2/(2ma^2)$,
the dimensionless eigenvalues and eigenvectors
are exactly
\begin{equation}
\boxed{\varepsilon_n(N)=4(N+1)^2
\sin^2\left(\frac{n\pi}{2(N+1)}\right)},\quad
\boxed{v^{(n)}_j=
\sqrt{\frac2{N+1}}\sin\frac{jn\pi}{N+1}},
\end{equation}
for $n,j=1,\ldots,N$.
This follows by applying the discrete second
difference to a sine eigenvector and using
$\sin(\theta+\delta)+\sin(\theta-\delta)
=2\sin\theta\cos\delta$.
For small sizes:
\begin{equation}
\begin{array}{c|l}
N & \{\varepsilon_n\}\\\hline
1& 8\\
2&9,\ 27\\
3&9.372583,\ 32,\ 54.627417
\end{array}
\end{equation}
For $N=10$ and $100$, the lowest five are:
\begin{center}
\begin{tabular}{rrrr}
\toprule
$n$ & $\varepsilon_n(10)$&$\varepsilon_n(100)$
& $\varepsilon_n(\infty)=n^2\pi^2$\\
\midrule
1&9.802700&9.868809&9.869604\\
2&38.416645&39.465687&39.478418\\
3&83.523702&88.762003&88.826440\\
4&141.469567&157.710064&157.913670\\
5&207.559809&246.243168&246.740110\\
\bottomrule
\end{tabular}
\end{center}
The convergence is of order $(\Delta x)^2$
for fixed low-lying $n$, and the eigenvectors
approach the corresponding continuum sine
functions.

\problem{2.62}{Numerical energies for a sinusoidal well bottom}
Let the infinite well occupy $0<x<a$
with $V(x)=500E_*\sin(\pi x/a)$,
$E_*=\hbar^2/(2ma^2)$.
For $N=100$ interior grid points
$x_j=ja/(N+1)$, use the dimensionless
tridiagonal matrix
\begin{equation}
\boxed{H_{jj}/E_*=2(N+1)^2+
500\sin\frac{\pi j}{N+1},\quad
H_{j,j\pm1}/E_*=-(N+1)^2}.
\end{equation}
Diagonalization yields the three lowest levels
(in units of $E_*$):
\begin{equation}
\boxed{\varepsilon_1\simeq304.65956,\quad
\varepsilon_2\simeq304.80037,\quad
\varepsilon_3\simeq476.16270}
\quad(N=100).
\end{equation}
The numerical eigenvectors are normalized by
$\sum_j|v_j|^2=1$ and can be converted to
continuum-normalized wavefunctions with
$\psi_j=v_j/\sqrt{\Delta x}$.
The first two states are close in energy
because the strong central positive hump
creates two effective edge wells with a
small tunneling splitting.
\check{The indicated $N=100$ matrix results
are specific discretization outputs, not
continuum-exact eigenvalues; the companion
script recomputes them and generates plots.}

\problem{2.63}{Oscillator partition function and Einstein heat capacity}
For equilibrium probabilities
$P_n=\ee^{-\beta E_n}/Z$ with
$Z=\sum_n\ee^{-\beta E_n}$,
\begin{equation}
-\frac{\partial}{\partial\beta}\ln Z
=-\frac1Z\frac{\partial Z}{\partial\beta}
=\frac1Z\sum_n E_n\ee^{-\beta E_n}
=\boxed{\bar E}.
\end{equation}
With oscillator energies $E_n=(n+1/2)\hbar\omega$,
the geometric series gives
\begin{equation}
\boxed{Z=\frac{\ee^{-\beta\hbar\omega/2}}
{1-\ee^{-\beta\hbar\omega}}}.
\end{equation}
Differentiating its logarithm yields
\begin{equation}
\boxed{\bar E=\frac{\hbar\omega}{2}
+\frac{\hbar\omega}{\ee^{\beta\hbar\omega}-1}
=\frac{\hbar\omega}{2}
\frac{1+\ee^{-\beta\hbar\omega}}
{1-\ee^{-\beta\hbar\omega}}}.
\end{equation}
For three oscillators per atom, define
$\theta_E=\hbar\omega/k_B$.
Then $C=3\,\partial\bar E/\partial T$ gives
\begin{equation}
\boxed{C=3k_B
\left(\frac{\theta_E}{T}\right)^2
\frac{\ee^{\theta_E/T}}
{(\ee^{\theta_E/T}-1)^2}}.
\end{equation}
At $T\gg\theta_E$,
$\ee^{\theta_E/T}-1\sim\theta_E/T$, hence
$C\to3k_B$ (classical Dulong--Petit limit).
At $T\ll\theta_E$,
$C\sim3k_B(\theta_E/T)^2\ee^{-\theta_E/T}\to0$.
The textbook's diamond curve illustrates
this temperature dependence qualitatively;
the model does not capture all real-crystal
phonon dispersions.

\problem{2.64}{Power-series solution of Legendre's equation}
For real $\ell\ge0$, start with
\begin{equation}
(1-x^2)f''-2xf'+\ell(\ell+1)f=0,\qquad
f(x)=\sum_{n=0}^{\infty}a_nx^n.
\end{equation}
Substituting
$f''=\sum_n(n+2)(n+1)a_{n+2}x^n$
and collecting the coefficient of $x^n$,
\begin{equation}
(n+2)(n+1)a_{n+2}
+[\ell(\ell+1)-n(n+1)]a_n=0.
\end{equation}
Therefore
\begin{equation}
\boxed{a_{n+2}=
\frac{n(n+1)-\ell(\ell+1)}
{(n+1)(n+2)}\,a_n}.
\end{equation}
Even and odd coefficients form independent
series. For integer $\ell$, the series
of matching parity truncates at $n=\ell$.
The resulting low-degree polynomial solutions,
normalized by $P_\ell(1)=1$, are
\begin{align}
\boxed{P_0(x)}&=\boxed{1},&
\boxed{P_1(x)}&=\boxed{x},\\
\boxed{P_2(x)}&=\boxed{\frac12(3x^2-1)},&
\boxed{P_3(x)}&=\boxed{\frac12(5x^3-3x)}.
\end{align}
For noninteger $\ell$, an isolated parity
series does not terminate and generically has
a logarithmic endpoint singularity. More general
linear combinations, however, include the
non-polynomial Legendre function $P_\ell(x)$,
which is regular at one endpoint. Imposing
regularity at \emph{both} endpoints in the usual
self-adjoint eigenvalue problem selects
nonnegative integer $\ell$.

\section{Numerical visualization and convergence tests}
The script \texttt{scripts/numerical\_checks.py} checks Problems
2.27, 2.48, 2.59, 2.61, and 2.62 independently.
The plots use the stated dimensionless units.
\IfFileExists{figures/bouncing_ball.pdf}{
\begin{figure}[htbp]\centering
\includegraphics[width=.80\linewidth]{figures/bouncing_ball.pdf}
\caption{Normalized Airy eigenstates for the quantum bouncing ball,
Problem 2.59.}
\end{figure}}{}
\IfFileExists{figures/box_convergence.pdf}{
\begin{figure}[htbp]\centering
\includegraphics[width=.80\linewidth]{figures/box_convergence.pdf}
\caption{Convergence of finite-difference energies toward
$n^2\pi^2$ for Problem 2.61.}
\end{figure}}{}
\IfFileExists{figures/sinusoidal_well.pdf}{
\begin{figure}[htbp]\centering
\includegraphics[width=.80\linewidth]{figures/sinusoidal_well.pdf}
\caption{First three eigenfunctions of the sinusoidal-bottom
well with $N=100$, Problem 2.62.}
\end{figure}}{}
\clearpage

\section*{References and scope}
\addcontentsline{toc}{section}{References and scope}
\begin{enumerate}[label={[\arabic*]},leftmargin=16pt]
\item D. J. Griffiths and D. F. Schroeter,
\textit{Introduction to Quantum Mechanics}, 3rd edition, Cambridge
University Press (2018), Chapter 2, textbook pages 25--90.
DOI: \href{https://doi.org/10.1017/9781316995433}{10.1017/9781316995433}.
\item GitHub project: \url{https://github.com/KrAzad0/QM-CH1-SOLUTION}.
\end{enumerate}
Problem statements have been paraphrased to avoid reproducing copyrighted
textbook questions. All mathematical work is independently written.
Numerical claims should be checked against the indicated computational
scripts, parameter choices, and boundary conditions.
\end{document}
